\documentclass[a4paper,12pt]{article}

\usepackage{alltt, fancyvrb, url}
\usepackage{graphicx}
\usepackage[utf8]{inputenc}
\usepackage{float}
\usepackage{hyperref}
\usepackage[italian]{babel}
\usepackage[table]{xcolor}
\usepackage{array}
\usepackage{multirow}
\usepackage{titlesec}
\usepackage{listings}
\usepackage[italian]{cleveref}
\usepackage{acronym}
\usepackage{subcaption}

\renewcommand{\thesection}{\arabic{section}}

\sloppy
\setlength{\parindent}{0pt}

\title{\textbf{Meetings}}
\author{}
\date{}

\begin{document}
\maketitle
\tableofcontents

\newpage

Al termine di ogni meeting il verbale stilato viene condiviso con tutti i partecipanti per la revisione e
l'approvazione. Il verbale comprende le decisioni prese ed i punti chiave discussi, oltre ad eventuali
promemoria per futuri incontri.

\section{Kick-off meeting}
Partecipanti: Project Manager del fornitore e del cliente (ChargeAndTrack), team del cliente, membri principali del
team di sviluppo.\\

Agenda:
\begin{itemize}
    \item Introduzione;
    \item presentazione dello sponsor;
    \item presentazione degli aspetti rilevanti del progetto da parte dello sponsor;
    \item presentazione del progetto da parte del cliente;
    \item presentazione del progetto da parte del project manager del fornitore;
    \item presentazione del team di sviluppo;
    \item discussione del PDS.
\end{itemize}

\section{Meeting interno}
Partecipanti: team di sviluppo, Project Manager del fornitore.\\

Agenda:
\begin{itemize}
    \item Introduzione;
    \item assegnazione responsabilità per i task;
    \item definizione delle regole operative del team;
    \item finalizzazione della schedula;
    \item identificazione e scrittura dei Work Packages.
\end{itemize}

\end{document}